\section{GPAT Overview}
The Gridded Plume Analysis Tool (GPAT) derives from work done by the University of Bristol Atmospheric Chemistry Research Group (ACRG), Professor Ulrich Schumann and colleagues at Deutsches Zentrum für Luft- und Raumfahrt (DLR), as well as Breakthrough Energy and their recent development of the \textit{pycontrails} open-source repository for aviation climate impact modelling \cite{pycontrails2024}. Each group's respective work was integral to the development and application of GPAT, as will become apparent throughout this chapter.

This tool was built out of a development version of the pycontrails library (v0.53.2), with the majority of the models and interlinking software being written in Python. The only exception to this being the photochemical box model (BOXM v2.0) written in FORTRAN 90, which is an adaptation of a legacy model (BOXM v1.0) \cite{Derwent1996} built in FORTRAN 77. All libraries, packages and tools used throughout the software development phase will be outlined explicitly where relevant.

The primary purpose of GPAT is to simulate a spatial domain in which to plot aircraft trajectories, estimate their emissions, model the dispersion and advection of the exhaust plumes carrying the emissions, then to compute the evolution of atmospheric chemistry and physics across the domain. The software architecture of the modelling framework is presented in figure~\ref{gpat_diag}. This portrays the input parameters, constituent models and the input and output datasets related to each model. 

\begin{landscape}
    \begin{figure}[htbp]
        \centering
        \hspace{-1.45cm}	
        \includegraphics[height=0.85\textwidth, width=1.387\textheight]{5/gpat_diag.pdf}
      
        \caption{GPAT process flow schematic portraying the input parameters (rhombuses), constituent models (boxes) and flow of data through a range of datasets (cylinders).}
        %\vspace{-0.5cm}
        \label{gpat_diag}
    \end{figure}
\end{landscape}
